% Answers to Chapter 8's exercises, lectures/L08/appendix/b_exercises.md. The book's own answers:
% the lecture publishes none.
\answerchapter{8}{The Watchdog Timer}

\answer{c8:ex:1}
Against the program that has gone wrong switching off the thing that is there to catch it going
wrong. A watchdog earns its keep exactly when the code is misbehaving: stuck in a loop, running off
through a corrupted pointer or return address, writing to registers it never meant to touch. If a
single stray write to \code{WDTCSR} could clear \code{WDE} or stretch the timeout, the same fault
the watchdog should recover from could disarm it first, and the chip would hang anyway.

With the interlock, one accidental write is not enough: turning the watchdog off or changing its
timeout takes \code{WDCE} and \code{WDE} written together, and then the new configuration written
with \code{WDCE} clear within four cycles. Runaway code is very unlikely to produce that exact
sequence by chance. (The datasheet is precise about which direction is protected: setting
\code{WDE} is allowed at any time, because a stray write that arms the watchdog does no harm; it is
clearing \code{WDE} and changing the prescaler that need the sequence.) EEPROM's \code{EEMPE}
guards stored data against an accidental write; the watchdog's \code{WDCE} guards the safety net
itself. \code{WDRF} (\secref{c8:a:3}) is the same idea taken one step further: after a watchdog
reset it holds \code{WDE} set until the flag is cleared.
